\chapter{Conclusions and Recommendations}\label{ch:conclusion}

\section{Introduction}\label{sec:concl-intro}

This chapter summarises the key findings of the dissertation, addresses each research question, states the contributions, acknowledges limitations, and presents recommendations for future work.

\section{Conclusions}\label{sec:concl-conclusions}

The NeuroShell IRE successfully implements a guarded, schema-constrained intent recognition pipeline that treats a local LLM as an unreliable proposition generator. The seven-stage pipeline (S1--S7 with S6B) combined with five middleware components (M1--M5) provides defence-in-depth through deterministic validation, adversarial detection, RBAC enforcement, session context, semantic caching and consistency checks.

Empirical evaluation (Exp1--Exp4) demonstrates that the guarded pipeline significantly outperforms raw LLM output. On the holdout set, Catch improved from 0\% to 72\% (Gemma) and 0\% to 64\% (Qwen), while FP reduced from ~32\% to ~14--15\% and DAcc increased by 8--13 percentage points. Both backends benefit from the guardrail stack, confirming generalisability (RQ3). Ablation analysis confirms that layered deterministic validation (schema/regex/network/scope) provides the largest joint contribution to flaw detection, and no single layer is sufficient (RQ2). Overall, the stacked approach is effective for improving both safety and usability (RQ1).

The implementation is fully evidence-based: all components exist in the codebase (\texttt{src/}), are configured via \texttt{config/}, and are validated by unit/integration tests. The system operates entirely offline with no retrieval at inference time, matching the design goals.

\section{Research Questions Addressed}\label{sec:concl-rq}

\begin{itemize}
  \item \textbf{RQ1:} The stacked validation pipeline improves DAcc by +8--13 pp, Catch by +64--72 pp, and reduces FP by -16--18 pp compared to raw LLM output.
  \item \textbf{RQ2:} The combined deterministic validation stack (S5 schema, S6 regex, S6B network, S7 scope) contributes most to catch rate; other layers (adversarial detection M1, RBAC M2/M5, self-consistency, semantic cache M4) provide complementary coverage.
  \item \textbf{RQ3:} Guardrails generalise across Gemma 4 e4b and Qwen2.5-Coder 7B under identical CPU-only configuration, with Gemma achieving higher absolute performance.
\end{itemize}

\section{Contributions}\label{sec:concl-contrib}

Key contributions of this work include: (1) implementation of a seven-stage guarded pipeline with five middleware components for offline intent recognition in a security context; (2) deterministic validation enforcing schema, entity formats, network architecture constraints, and scope policy; (3) defence-in-depth via three-layer adversarial detection and pre/post RBAC enforcement; (4) empirical evaluation using a 900-record dataset with ablation across guardrail layers and backend comparison; (5) evidence-based documentation grounded strictly in the implemented codebase.

\section{Limitations}\label{sec:concl-limitations}

Results are single-run point estimates without statistical confidence intervals. Dataset category sizes are non-uniform. Evaluation was conducted on specific hardware and model versions (CPU-only). The threat model excludes physical compromise and indirect injection via retrieved content (by design, no retrieval). These limitations are acknowledged and qualify the reported results.

\section{Recommendations for Future Work}\label{sec:concl-recommendations}

Future work should include: (1) repeated runs with multiple seeds and reporting of confidence intervals/standard deviations; (2) expansion and balancing of the evaluation dataset across flaw categories; (3) evaluation on additional hardware configurations and newer model versions; (4) exploration of lightweight local verification methods to further reduce FP while maintaining high Catch; (5) integration testing with downstream planner/executor components under controlled laboratory conditions; (6) systematic red-teaming against the implemented guards to identify residual attack vectors.
